\section{Discussion}
\label{sec:discussion}

\para{What the evidence supports.} The strongest supported claim is that \methodname is a useful triage workflow. It connects four pieces of evidence that are usually separated: residue-level sparse activations, \dms sensitivity, known-motif exclusion, and internal model intervention. The workflow surfaces candidates that are plausible enough for follow-up while making the uncertainty visible.

\para{What the evidence does not support.} The experiment does not establish a new biological mechanism. The aggregate association results are close across \sae, \rawhidden, and \shuffledsae controls. This matters because a sparse feature can be interpretable without being a distinct biological abstraction. It also matters because 10,240 sparse features tested on sequences of length 37--52 create many chances to find high correlations. In this setting, raw hidden dimensions and shuffled activations are not weak baselines; they are necessary checks against overclaiming.

\para{Interpreting feature 6419.} Feature 6419 is interesting because it passes the implemented motif filters and its ablation changes \esm masked-marginal scores more at high-activation positions. This suggests that the feature participates in the model's mutation-score computation for the \texttt{YNZC\_BACSU} target. The next question is biological coherence. High-activation positions should be tested for structural proximity, solvent exposure, conservation, physicochemical patterning, and consistency across homologs. Without those checks, the feature remains a model-internal hypothesis rather than an inferred mechanism.

\para{Limitations.} The pilot uses \esmeight for tractability, while larger \esm models may contain cleaner or different sparse features. The assay panel contains 16 short proteins, mostly from Tsuboyama-style stability assays, so it is not broad biological coverage. The \prosite converter covers common pattern syntax but not all profiles or rules. The \elm/\prosite scans are permissive and incomplete, so ``low motif alignment'' means low alignment under the implemented scan, not absence of known biology. The feature-level permutation budget is small relative to 10,240 tested features, so candidate $q$-values do not survive FDR correction. Finally, internal ablation effects are measured in masked-marginal log-probability units, not experimental fitness.

\para{Broader implications.} A cautious workflow is valuable because PLM-derived biological claims can otherwise move faster than the evidence. Feature triage can help prioritize structural analysis and mutagenesis, but it should not replace annotation curation or experimental validation. The same caution also reduces misuse: model-internal signals should not be treated as sufficient evidence for functional annotation, disease interpretation, or protein-design decisions without independent checks.
